% !TeX root = ../thesis.tex
% Vorläufige Gliederung mit geplanten Inhalten.
\chapter{Ergebnisse und Analyse}
\label{chap:ergebnisse}
\section{Ergebnisse der Extraktionsqualität}
\begin{inhalte}
  \item Vergleich der untersuchten Ansätze
  \item Auswertung der zentralen Qualitätsmetriken
\end{inhalte}
\section{Analyse der untersuchten Ansätze}
\begin{inhalte}
  \item Gegenüberstellung der untersuchten Modellklassen und Verarbeitungspipelines
  \item Unterschiede zwischen spezialisierten Modellen, Vision-Language-Modellen und hybriden Ansätzen
\end{inhalte}
\section{Fehleranalyse}
\begin{inhalte}
  \item Typische Fehlerarten
  \item Unterschiede zwischen Modellklassen
  \item Einfluss von Handschrift und Dokumentqualität
\end{inhalte}
\section{Analyse der Effizienz}
\begin{inhalte}
  \item Laufzeit
  \item Rechenaufwand
  \item Anzahl beziehungsweise Komplexität der Inferenzschritte
\end{inhalte}
\section{Kostenanalyse}
\begin{inhalte}
  \item \acs{api}-Kosten
  \item Lokale Inferenz
  \item Kosten pro Dokument beziehungsweise Extraktion
  \item Manueller Aufwand für Vorbereitung, Kontrolle und Korrektur
  \item Vergleich mit der bisherigen manuellen Erfassung bei gleichem Bearbeitungsstand
  \item Einmalige Aufwände und bedingte Szenarien getrennt ausweisen
\end{inhalte}
\section{Analyse der praktischen Einsatzfähigkeit}
\begin{inhalte}
  \item Technische Umsetzbarkeit
  \item Integrationsaufwand
  \item Hardware- und Infrastrukturbedarf
  \item Datenschutz und Informationssicherheit
  \item Skalierbarkeit und Automatisierbarkeit im Unternehmenskontext
\end{inhalte}
\section{Gesamtbewertung der untersuchten Ansätze}
\begin{inhalte}
  \item Gegenüberstellung der Ergebnisse zu Qualität, Effizienz, Kosten und praktischer Einsatzfähigkeit
\end{inhalte}
